\chapter{Présentation de l'entreprise et du contexte du stage}

Ce chapitre présente l'entreprise d'accueil, Mövenpick Hôtel de Tanger, son secteur
d'activité, son historique, son organisation interne ainsi que son environnement
(clientèle, fournisseurs, concurrence). Il présente également le département
Informatique (IT) au sein duquel le stage s'est déroulé, son rôle dans l'organisation de
l'hôtel, ainsi que la mission confiée à l'étudiant durant ces deux mois.

\section{Présentation de l'entreprise}

\subsection{Activité de l'entreprise}

Le Mövenpick Hôtel de Tanger fait partie du groupe Mövenpick Hotels \& Resorts, une
chaîne hôtelière internationale opérant dans le secteur de l'hôtellerie et du tourisme.
L'établissement propose des prestations d'hébergement, de restauration et
d'organisation d'événements, destinées principalement à une clientèle d'affaires et de
loisirs.

\subsection{Historique de l'entreprise}

Le groupe Mövenpick a été fondé en 1948 à Zurich (Suisse) par Ueli Prager, initialement
sous la forme d'un concept de restauration. Son activité hôtelière a débuté en 1973 avec
l'ouverture de ses deux premiers hôtels à Zurich. Depuis septembre 2018, Mövenpick
Hotels \& Resorts est intégralement détenu par le groupe Accor, ce qui lui permet de
bénéficier de son réseau de distribution et de ses outils opérationnels à l'échelle
internationale. Le groupe est aujourd'hui présent dans plusieurs dizaines de pays, avec
une forte présence en Europe et au Moyen-Orient.

\subsection{Organisation interne de l'entreprise}

L'hôtel est structuré en plusieurs départements assurant chacun une fonction
spécifique : la Direction Générale, les Ressources Humaines, la Finance, l'Hébergement
(Front Office, Housekeeping), la Restauration, ainsi que le département Informatique
(IT), au sein duquel s'est déroulé ce stage. Ce dernier est chargé de la gestion du
réseau informatique interne et du support technique aux différents services de l'hôtel,
tandis que les systèmes métiers principaux (tels que le Property Management System)
sont gérés par un prestataire externe spécialisé.

\section{Environnement de l'entreprise}

\subsection{Clientèle}

La clientèle du Mövenpick Hôtel de Tanger se compose principalement de voyageurs
d'affaires, de touristes internationaux et de participants à des événements ou
séminaires organisés au sein de l'établissement. Tanger, en tant que ville portuaire et
touristique, attire une clientèle variée tout au long de l'année.

\subsection{Fournisseurs}

L'hôtel fait appel à différents fournisseurs et prestataires externes, notamment pour la
restauration, la maintenance des équipements, et la gestion des systèmes informatiques
métiers (PMS), assurée par un prestataire spécialisé externe au département IT interne.

\subsection{Concurrence}

Le marché hôtelier de Tanger est marqué par la présence de plusieurs établissements
internationaux et locaux, notamment dans les segments haut de gamme et affaires. Le
Mövenpick Hôtel se positionne comme un établissement haut de gamme, misant sur la
qualité de service et son emplacement stratégique dans la ville.

\bigskip
La seconde partie de ce chapitre présente plus en détail le département Informatique
(IT), son rôle au sein de l'hôtel, ainsi que la mission confiée à l'étudiant durant ce
stage.

\section{Aperçu sur le département d'attache}

\subsection{Missions du département}

Le département Informatique assure la gestion et la maintenance du système
d'information de l'hôtel : réseau informatique, postes de travail, imprimantes et
équipements divers. Les systèmes métiers principaux (PMS, POS) sont gérés par un
prestataire externe spécialisé, ce qui recentre l'activité du département IT interne sur
l'infrastructure réseau et le support technique.

\subsection{Lien fonctionnel avec le reste de l'entreprise}

Le département IT intervient de façon transversale : il fournit un support technique à
l'ensemble des autres départements de l'hôtel (Front Office, Restauration, Ressources
Humaines, etc.) et assure la disponibilité et la sécurité du réseau informatique utilisé
par l'ensemble du personnel.

\subsection{Métiers du département}

Le département IT est composé de techniciens et responsables informatiques chargés du
support utilisateur, de la maintenance du réseau et des équipements, ainsi que du suivi
des prestataires externes pour les systèmes métiers. Le travail y est à la fois
technique (réseau, matériel) et relationnel (support aux autres départements).

\section{Description de la mission}

\subsection{Position de l'étudiant}

Au sein du département IT, l'étudiant a été chargé de la conception et du développement
d'un outil interne de découverte et de cartographie des équipements connectés au réseau
informatique de l'hôtel, en complément des tâches de support courantes du département.

\subsection{Tâches réalisées}

Les principales tâches confiées portent sur l'analyse des besoins du département en
matière de suivi réseau, la mise à niveau sur les concepts réseau nécessaires, ainsi que
la conception, le développement et le déploiement de l'outil décrit dans le chapitre
suivant.

\subsection{Planification des tâches}

La planification prévisionnelle du stage, détaillée dans le cahier des charges du
projet, s'étend sur les deux mois du stage : mise à niveau réseau, développement du
backend et de la base de données, développement du tableau de bord, implémentation de la
détection d'anomalies, puis conteneurisation, déploiement et documentation.
